\documentclass[letterpaper,12pt]{letter}


\usepackage{geometry}
\usepackage{graphicx}
\usepackage{fancyhdr}
\usepackage{marvosym}
\usepackage{hyperref} 

\usepackage{lmodern}
\usepackage{cancel}
\usepackage{graphicx}
\usepackage{braket}
\usepackage{latexsym,amssymb,amsmath}
\usepackage{pdfpages}
\usepackage{xcolor}
\usepackage{capt-of}
\usepackage{amsmath}
\usepackage{cite}

\renewcommand{\headrulewidth}{0pt}

\fancypagestyle{plain}{
\fancyhf{}
\fancyhead[L]{\includegraphics[width=0.6\textwidth]{logo.png}}
\fancyfoot[C]{\Email~{\normalfont\ttfamily mulhearn@physics.ucdavis.edu}}
}

\pagestyle{plain}

\makeatletter
\let\ps@empty\ps@plain
\let\ps@firstpage\ps@plain
\makeatother

\address{Michael Mulhearn \\ One Shields Ave \\ Davis, CA 95616}

\addtolength{\oddsidemargin}{-0.25in}
\addtolength{\evensidemargin}{-0.25in}


\def\letterparta{
We are required by the Undergraduate Council to consult your
department~[1] regarding proposed changes to our physics undergraduate major requirements.  Despite the fact that {\bf we expect the
  impact on your department to be minimal}, we will explain it here
fully.  We hereby request (A) a {\bf vote by your department} on
whether or not to support our changes to the physics major
requirements, and (B) a {\bf letter of support} from your department
that includes the outcome of that vote.

Our proposal~[2] has several aims, the most important of which I
consider to be: (1) improving the experience for students that
transfer to UCD from a community college in their junior year, (2)
increasing the role of computational physics in our program, (3)
improving our laboratory offerings, and (4) providing more schedule
flexibility and elective choices for our students.

The full proposal is quite long.  Therefore, this letter includes a
much shorter summary of the courses taught in your department which
are impacted by this proposal, and our best estimate of the size of
that impact.
}

\def\letterpartb{ The approval process for physics has proven long and
windy, with many unexpected delays.  I would be happy to share my
insights with you some day if your department decides to update your
own program!  As we enter year five of this process for me personally,
I am pushing hard to return the proposal to the Undergraduate Council,
including your department's letter, as early in the spring quarter as
possible. These required consultations are the last item that we need,
and {\bf I would be much obliged for your department's rapid response},
which would enable our department to finish our proposal this year.

        \closing{Sincerely,\\
        \fromsig{\includegraphics[height=0.3in]{github-signature.png}} \\
        \fromname{Michael Mulhearn \\ Professor \\ Department of Physics}}

{\bf References}\\[5pt]
[1] The UCD requirements for program revisions, including this required consultation:
\begin{quote}
 \url{https://ucdavis.box.com/s/v1682jq4dp9sgcrhyr3oog1wp3pa36fd}
\end{quote}
[2] The latest version of the proposed update to the undergraduate physics curriculum:
\begin{quote}
 \url{https://github.com/mulhearn/classwork/blob/master/curriculum/curriculum.pdf}
\end{quote}
}




\begin{document}

%----------------------------------------------------------------------------------
%  Mathematics, Chair Anne Schilling (MAT) 
%----------------------------------------------------------------------------------

\begin{letter}{Department of Mathematics \\ anne@math.ucdavis.edu}

\opening{Dear Chair Anne Schilling,}

\letterparta

{\bf Changes for Physics and Applied Physics Majors:} All of our
physics majors (BS in Physics, BS in Applied Physics, and AB in
Physics) share a common set of math requirements that cover calculus,
linear algebra, and differential equations.  As part of the upgrade of
our undergraduate physics curriculum, we have updated the math
requirements for all physics majors to include additional options now
offered by the math department, specifically linear algebra and
differential equations with applications to biology (MAT 27A and MAT
27B) and Modern Linear Algebra (MAT 67).  The new requirements are
listed in Table 1 (below).

{\bf Expected Impact:} In principle, these changes impact all of our
250 physics majors, with an average annual throughput in any math
class of approximately 50.  However, the actual impact will be much
smaller.  We have already been accepting these courses as equivalent
on a case by case basis whenever requested by a student.  These
changes simply advertise this equivalence.  We expect that this change
will have very little impact: even if a surprising large $10\%$ of
physics majors decide to use an alternative class, the impact to the
math department would be about five students per year per course.

We also understand this update to be fully consistent with the math
departments intentions in offering these alternative courses.  Please
note that {\bf the inclusion of MAT 27A and 27B was specifically
  requested by the math department}, as shown in the enclosed letter.

{\bf Table 1}: Preparatory Subject Matter (MAT)\\
\begin{tabular}{|llllll|}
\hline
Course & & Units & Offered & Prereqs & Name \\
\hline
MAT & 21A & 4 & FWS & & Differential Calculus\\ 
    & 21B & 4 & FWS &  & Integral Calculus \\ 
    & 21C & 4 & FWS &  & Partial Derivatives and Series\\ 
    & 21D & 4 & FWS &  & Vector Analysis\\
\hline
    & one of:  & & & & \\
MAT & 22A & 3 & FWS &  & Linear Algebra\\
    & 27A & 3 & FWS &  & Linear Algebra\\
    & 67  & 4 & FWS &  & Mod. Linear Algebra\\
\hline
MAT & 22B & 3 & FWS &  & Differential Equations\\
    & or  & & & & \\
MAT & 27B & 3 & FWS &  & Differential Equations\\ 
\hline
\end{tabular}\\


{\bf Update to Math prerequisites:} As part of our curriculum
overhaul, we have also updated all of the math prerequisites to be
consistent with the equivalencies used by the math department, as
shown in the Table 2 (on the next page).

As physics majors are required to take the MAT 21 sequence, these
changes only impact non-physics majors taking the PHY 9 sequence.
Therefore, {\bf these changes are outside the scope of your department
  vote on the proposed changes to physics majors}.  They have been
included here only for completeness, and we certainly welcome any
comments on the appropriateness of these changes.

\newpage

{\bf Table 2:}  Math prerequisite for lower division courses related to physics
courses.  Where indicated, course grades are the minimum accepted to
meet prerequisites.  {\bf The prerequisites for the math courses are
  set by the math department and are only reported here}, including
some simplifications.  Where possible, math prerequisites for physics
courses will adopt the same choices made by the math department with
respect to equivalent math coursework.
\begin{center}
\begin{tabular}{|llll|}
\hline
Course & & Math Prereqs (Minimum Grade)& Name \\
\hline
MAT & 16A & & Short Calc.\\
    & 16B & 21A/21AH/17A/16A(C-) & Short Calc.\\
    & 16C & 21B/21BH/17B/16B(C-) & Short Calc.\\
    & 17A &                      & Calc (Bio\&Med)\\
    & 17B & 21A/21AH/17A/16A(C-) & Calc (Bio\&Med)\\
    & 17C & 17B(C-)              & Calc (Bio\&Med)\\
    & 21A &                      & Calculus \\
    & 21B & 21A/21AH(C-) or 17A(B) & Calculus \\
    & 21C & 21B/21BH/17C/16C(C-) or 17B(B) & Calculus\\
    & 21D & 21C/21CH(C-) or 17C(B)  & Calculus \\
    & 22A & 21C/21CH/17C/16C(C-)    & Linear Algebra \\
    & 22B & 67/22A(C-)              & Diff. Eqs. \\
    & 27A & 21C/21CH/17C(C-)        & Linear Algebra w/ App. Bio. \\
    & 27B & 27A/22A(C-)             & Diff. Eqs. w/ App. Bio.\\
    & 67  & 21C/21CH(C-)            & Mod. Lin. Alg. \\
\hline
\hline
PHY & 9A & 21B/21BH/17C/16C(C-) or 17B(B) & Class. Physics\\
    & 9B & 21C/21CH(C-) or 17C(B) & Class. Physics\\
    & 9C & 21D(C-) & Class. Physics\\
    & 9D & 22A(C-) or 27A(C-) or 67(C-) & Class. Physics\\
    &    & $\parallel$22B or $\parallel$27B & \\
\hline
\hline
PHY & 9HA & $\parallel$21B/$\parallel$21BH & Hon. Physics\\
    & 9HB    & 21B/21BH(C-) & Hon. Physics \\
    & 9HC    & 21C/21CH & Hon. Physics \\
    & 9HD    & 21D  & Hon. Physics \\
\hline
PHY & 45     & $\parallel$22B or $\parallel$27B & Computational Physics \\
\hline
\end{tabular}
\end{center}
\newpage

\letterpartb
        
\end{letter}

%----------------------------------------------------------------------------------
%  Earth and Planetary Sciences, Chair Kari Cooper (GEL) 
%----------------------------------------------------------------------------------

\begin{letter}{Department of Earth and  Planetary Science \\ kmcooper@ucdavis.edu}
  
\opening{Dear Chair Kari Cooper,}

\letterparta

{\bf No changes to GEL courses for Physics AB and BS majors:}
There are no changes to Geology coursework to the Physics majors (AB
and BS) including the Astrophysics specialization.  The proposed
requirements for the Astrophysics specializations continues to include
GEL 163 as an elective course, as in the current requirements.

{\bf Changes to GEL courses for Applied Physics majors:} The proposed
physics curriculum includes changes to the Applied Physics major with
specializations in Atmospheric Physics, Physical Oceanography, and
Geological Physics related to GEL coursework, as detailed in
Tables~1-3.  All three of these majors were designed before
significant changes to the prerequisite structure of GEL courses.

In the current requirements for the Atmospheric physics
specialization, GEL 150A is a required course, but due to the
extensive prerequisites, this is no longer feasible while remaining
below the unit cap on required coursework.  GEL 116N is a possible
elective.  In the proposed requirements, GEL 50 is now a required
course (was only explicitly recommended, though implicitly required).
In addition, GEL 116N, 150A, and 150B are listed along with PHY and
ATM coursework as possible electives.

In the current requirements for the Physical Oceanography
specialization, GEL 116N and 150A were required. The extensive
prerequisites for these courses push the major past the unit cap
guidelines for required coursework.  In the proposed major, GEL 50,
55, 116N, and 150A are now all explicitly required courses.  Students
may also include 150B as one of two electives that include other
choices from PHY and ATM courses.

In the current requirements for the Geological Physics major, GEL 161
and 162 are required, as is one of GEL 146 or 163.  Due to
prerequisites, this now implicitly requires GEL 50 as well, which
pushes this major past the unit cap guidelines for required
coursework.  In the proposed major, GEL 50 and 50L are now explicitly
required courses.  Students can choose five electives from GEL 116N,
146, 150A, 150B, 161, 162, 163, and other PHY and ATM courses.  They
are required to select at least two from GEL 161,162, and 163.

To summarize, all three of these specializations have been changed to
properly account for prerequisite coursework for required courses and
to afford students additional choices for elective coursework in PHY,
ATM, and GEL.  

{\bf Expected impact of these changes on the Earth and Planetary Sciences department:}
Based on past enrollment, we expect these changes to impact at most
five students per year per course:
\begin{quote}
  {\bf GEL 50/50L/55/116N/146/150A/150B/161/162/163.}
\end{quote}
Furthermore, some of the changes simply make implicit requirements
explicit.  Also, these majors include a Program Variance, which
already allows students to substitute related coursework with approval
from the department.  So even on the current requirements, for
example, a student would easily receive permission to take GEL 150A as
an elective.  {\bf We therefore expect the impact of these changes on
  GEL course enrollment to be quite minimal.}

\newpage

{\bf Table 1:}  Required GEL Coursework for Atmospheric Physics
\vskip 0.25cm
\noindent
*: recommended, $^\#$: not offered every year, $\parallel$: concurrently.\\
\begin{tabular}{|llllll|}
\hline
Course & & Units & Offered & Prereqs & Name \\
\hline
GEL & 50     & 3 & FWS & & Physical Geology \\
\hline
\hline
    & Choose two of & & & & \\
\hline
PHY  & 105B   & 4 & & & \\
     & 105C   & 4 & $^\#$  &             & Continuum Mechanics\\
GEL  & 116N   & 3 & S  & GEL 50      & Oceanography\\
     & 150A   & 4 & S$^\#$ & GEL 116N,55 & Phys. and Chem. Oceanography\\
     & 150B   & 3 & W  & GEL 50      & Geological Oceanography\\
ATM  & 124    & 3 & F  & ATM 60      & Meteorological Inst. \& Obs. \\
     & 128    & 4 & W  & ATM 60      & Radiation and Satellite Met. \\
     & 158    & 4 & S  & ATM 121A    & Boundary-Layer Meteorology \\
\hline
\end{tabular}\\

{\bf Table 2:}  Required GEL coursework for Physical Oceanography
\vskip 0.25cm
\noindent
*: recommended, $^\#$: not offered every year, $\parallel$: concurrently.\\
\begin{tabular}{|llllll|}
\hline
Course & & Units & Offered & Prereqs & Name \\
\hline
GEL  & 50     & 3 & FWS & & Physical Geology \\
     & 55     & 3 & F      & CHE 2A & Intro. to Geochemistry\\
\hline
\hline
Course & & Units & Offered & Prereqs & Name \\
\hline
GEL  & 116N   & 3 & S      & GEL 50          & Oceanography\\
     & 150A   & 4 & S$^\#$ & GEL 116N,55 & Phys. and Chem. Oceanography\\
\hline
\hline
    & Choose two of & & & & \\
\hline
PHY  & 105B  & 4 & & & \\
     & 105C  & 4 & S$^\#$ & & Continuum Mechanics \\
ATM  & 115   & 3 & S   & ATM 60 & Hydroclimatology \\
     & 116   & 3 & F   & & Modern Climate Change \\
     & 120   & 4 & F   & ATM 60$\parallel$ & Atmos. Thermo. and Cloud Phys. \\
GEL  & 150B  & 3 & W & GEL 50 & Geological Oceanography\\
\hline
\end{tabular}\\

\newpage

{\bf Table 3:}  Required GEL coursework for Geological Physics
\vskip 0.25cm
\noindent
*: recommended, $^\#$: not offered every year, $\parallel$: concurrently.\\
\begin{tabular}{|llllll|}
\hline
Course & & Units & Offered & Prereqs & Name \\
\hline
GEL & 50     & 3 & FWS & & Physical Geology \\
    & 50L    & 2 &     & & Physical Geology Laboratory\\
\hline
\hline
    & Choose five of & & & & \\
\hline
PHY  & 105B   & 4 & &  & \\
     & 105C   & 4 & &  & Continuum Mechanics\\
ATM  & 120    & 4 & F  & ATM 60    & Atmos. Thermo. \& Cloud Phys. \\
     & 121A   & 4 & F  & ATM 120   & Atmospheric Dynamics \\
     & 121B   & 4 & W  & ATM 121A  & Atmospheric Dynamics \\
GEL  & 116N   & 3 &    & & Oceanography\\
     & 146    & 3 &    & & Rad. Iso. Geochem. \& Cosmochem. \\
     & 150A   & 4 & S$^\#$ & GEL 116N,55 & Phys. and Chem. Oceanography\\
     & 150B   & 3 & W  & GEL 50      & Geological Oceanography\\
\hline
     & incl at least two:  & & & & \\
\hline
GEL  & 161    & 3 & S$^\#$ & GEL 50 & Geophysical Field Methods \\
     & 162    & 3 & W$^\#$ & GEL 50 & Geophysics of the Solid Earth \\
     & 163    & 3 & F$^\#$ & GEL 50 & Planetary Geology \\
\hline
\end{tabular}\\

To summarize, all three of these specializations have been changed to
properly account for prerequisite coursework for required courses and
to afford students additional choices for elective coursework in PHY,
ATM, and GEL.

\letterpartb
        
\end{letter}

%----------------------------------------------------------------------------------
%  Land, Air \& Water Resources, Chair William Horwath (ATM)
%----------------------------------------------------------------------------------

\begin{letter}{Department of Land, Air and Water Resources \\ wrhorwath@ucdavis.edu}
  
\opening{Dear Chair William Horwath,}

\letterparta

{\bf No changes to ATM courses for Physics AB and BS majors:} There
are no changes to Geology coursework to the Physics majors (AB and
BS).

{\bf Changes to ATM courses for Applied Physics majors:} The proposed
physics curriculum includes changes to the Applied Physics major with
specializations in Atmospheric Physics, Physical Oceanography, and
Geological Physics related to ATM coursework, as detailed in
Tables~1-3.

In the current requirements for the Atmospheric Physics
specialization, ATM 120,121A, and 121B are required courses.  Due to
prerequisites, this now implicitly requires ATM 60 as well, which
pushes this major past the unit cap for required coursework.  In the
proposed requirements, ATM 60,120,121A and 121B are now all explicitly
required while the major still remains below the unit cap.  In the
current requirements, ATM 128 is a possible elective.  In the proposed
requirements, ATM 124,128, and 158 are included amongst PHY and GEL
courses as possible electives.

In the current requirements for the Physical Oceanography
specialization, ATM 120 is a required course.  This implicitly
requires ATM 60, which pushes this major past the unit cap for
required coursework.  In the proposed requirements, ATM 60 is a
recommended course, and ATM 115,116, and 120 are included as possible
electives alongside PHY and GEL courses.

In the current requirements for the Geophysics major, students must
choose one course from ATM 120,121A, and 121B.  Due to prerequisites,
this effectively requires students to take both ATM 60 and 120, 
pushing this major past the unit cap for required coursework.  In the
proposed requirements, ATM 60 is required and students choose five
electives from ATM 120,121A, 121B, and other courses in GEL and PHY.

To summarize, all three of these specializations have been changed to
properly account for prerequisite coursework for required courses and
to afford students additional choices for elective coursework in PHY,
ATM, and GEL.

{\bf Expected impact of these changes on the Land, Air \& Water Resources department:}
Based on past enrollment, we expect these changes to impact at most
five students per year per course:
\begin{quote}
{\bf ATM 60/115/116/120/121A/121B/124/128/158}.
\end{quote}
Furthermore, some of the changes simply make implicit requirements
explicit.  Also, these majors include a Program Variance, which
already allows students to substitute related coursework with approval
from the department.  So even on the current requirements, for
example, a student would easily receive permission to take ATM 158 as
an elective.
{\bf We therefore expect the impact of these changes on ATM course
  enrollment to be quite minimal.}

\newpage
{\bf Table 1:}  Required ATM coursework for Atmospheric Physics
\vskip 0.25cm
\noindent
*: recommended, $^\#$: not offered every year, $\parallel$: concurrently.\\
\begin{tabular}{|llllll|}
\hline
Course & & Units & Offered & Prereqs & Name \\
\hline
ATM & 60     & 4 & F   & & Intro. to Atmospheric Sci. \\
    & 120    & 4 & F   & ATM 60$\parallel$ & Atmos. Thermo. and Cloud Phys. \\
    & 121A   & 4 & W   & ATM 120 & Atmospheric Dynamics \\
    & 121B   & 4 & S   & ATM 121A & Atmospheric Dynamics \\
\hline
\hline
    & Choose two of & & & & \\
\hline
PHY  & 105B   & 4 & & & \\
     & 105C   & 4 & $^\#$  &             & Continuum Mechanics\\
GEL  & 116N   & 3 & S  & GEL 50      & Oceanography\\
     & 150A   & 4 & S$^\#$ & GEL 116N,55 & Phys. and Chem. Oceanography\\
     & 150B   & 3 & W  & GEL 50      & Geological Oceanography\\
ATM  & 124    & 3 & F  & ATM 60      & Meteorological Instr. \& Obs. \\
     & 128    & 4 & W  & ATM 60      & Radiation and Satellite Met. \\
     & 158    & 4 & S  & ATM 121A    & Boundary-Layer Meteorology \\
\hline
\end{tabular}\\

{\bf Table 2:}  Required ATM coursework for Physical Oceanography
\vskip 0.25cm
\noindent
*: recommended, $^\#$: not offered every year, $\parallel$: concurrently.\\
\begin{tabular}{|llllll|}
\hline
Course & & Units & Offered & Prereqs & Name \\
\hline
ATM  & 60*    & 4 & F   & & Intro. to Atmospheric Sci. \\
\hline
\hline
    & Choose two of & & & & \\
\hline
PHY  & 105B  & 4 & & & \\
     & 105C  & 4 & S$^\#$ & & Continuum Mechanics \\
ATM  & 115   & 3 & S   & ATM 60 & Hydroclimatology \\
     & 116   & 3 & F   & & Modern Climate Change \\
     & 120   & 4 & F   & ATM 60$\parallel$ & Atmos. Thermo. and Cloud Physics \\
GEL  & 150B  & 3 & W & GEL 50 & Geological Oceanography\\
\hline
\end{tabular}\\

\newpage
{\bf Table 3:}  Required ATM coursework Geological Physics
\vskip 0.25cm
\noindent
*: recommended, $^\#$: not offered every year, $\parallel$: concurrently.\\
\begin{tabular}{|llllll|}
\hline
Course & & Units & Offered & Prereqs & Name \\
\hline
ATM & 60     & 4 & F   & & Intro. to Atmospheric Sci. \\
\hline
\hline
    & Choose five of & & & & \\
\hline
PHY  & 105B   & 4 & &  & \\
     & 105C   & 4 & &  & Continuum Mechanics\\
ATM  & 120    & 4 & F  & ATM 60    & Atmospheric Thermo. \& Cloud Phys. \\
     & 121A   & 4 & F  & ATM 120   & Atmospheric Dynamics \\
     & 121B   & 4 & W  & ATM 121A  & Atmospheric Dynamics \\
GEL  & 116N   & 3 &    & & Oceanography\\
     & 146    & 3 &    & & Rad. Iso. Geochem. \& Cosmochem. \\
     & 150A   & 4 & S$^\#$ & GEL 116N,55 & Phys. and Chem. Oceanography\\
     & 150B   & 3 & W  & GEL 50      & Geological Oceanography\\
\hline
     & incl at least two:  & & & & \\
\hline
GEL  & 161    & 3 & S$^\#$ & GEL 50 & Geophysical Field Methods \\
     & 162    & 3 & W$^\#$ & GEL 50 & Geophysics of the Solid Earth \\
     & 163    & 3 & F$^\#$ & GEL 50 & Planetary Geology \\
\hline
\end{tabular}\\

\newpage
\letterpartb
        
\end{letter}

%----------------------------------------------------------------------------------
%  Chemistry, Chair David Goodin (CHE)
%----------------------------------------------------------------------------------
\begin{letter}{Department of Chemistry \\ goodinchair@ucdavis.edu}

\opening{Dear Chair David Goodin,}

\letterparta

{\bf No changes to CHE courses for Physics AB and BS majors:} There
are no changes to Chemistry coursework to the Physics majors (AB and
BS).

{\bf Changes to CHE courses for Applied Physics majors:} The proposed
physics curriculum includes changes to the Applied Physics major with
specializations in Physical Oceanography and Chemical Physics, as
detailed in Tables 1 and 2 (below).

In the proposed requirements for the Physical Oceanography
specialization, CHE 2A is now a required course.  This is to satisfy
the prerequisite for GEL 55 which is also a required course for this
major.

The current requirements for the Chemical Physics specialization
include CHE 2A,2B,2C and 124A as required courses.  There are no
chemistry electives, only a long list of recommended upper division
chemistry courses.  This is because both Chemistry and Physics are
highly hierarchical, and it is time consuming to complete the
prerequisites for upper division coursework in both subjects.  In the
proposed requirements, CHE 2A,2B,2C and 124A are still required.
However, students are also given the flexibility to take the sequence
CHE 110A,110B,110C, and 128A in the place of related coursework in
physics.  In additional, they may satisfy their lab requirements with
coursework in either physics or CHE 105,115,124L,129A,129B,129C.

These changes to the CHE requirements for the Chemical Physics
specialization are intended to encourage students to explore more
chemistry coursework, making this a more compelling and interesting
major.

{\bf Expected impact of these changes on the Chemistry department:}
Based on past enrollment, we expect these changes to impact at most
two students per year per course:
\begin{quote}
{\bf CHE 2A/105/110A/110B/110C/115/124L/128A/128B/129A/129B/129C}.
\end{quote}
Because CHE 2A was already implicitly required for the Physical
Oceanography specialization, making this requirement explicit should
have a negligible impact on the Chemistry department.

Also, these majors include a Program Variance, which already allows
students to substitute related coursework with approval from the
department.  So even on the current requirements, for example, a
student would likely receive permission to take CHE 110AB instead of
PHY 115AB.  {\bf We therefore expect the impact of these changes on
  CHE course enrollment to be quite minimal.}

\newpage

{\bf Table 1:}  Required CHE coursework for Physical Oceanography
\vskip 0.25cm
\noindent
*: recommended, $^\#$: not offered every year, $\parallel$: concurrently.\\
\begin{tabular}{|llllll|}
\hline
Course & & Units & Offered & Prereqs & Name \\
\hline
CHE  & 2A     & 5 & FW  & & \\
\hline
\end{tabular}\\

{\bf Table 2:}  Required CHE coursework for Chemical Physics
\vskip 0.25cm
\noindent
*: recommended, $^\#$: not offered every year, $\parallel$: concurrently.\\
\begin{tabular}{|llllll|}
\hline
Course & & Units & Offered & Prereqs & Name \\
\hline
\hline
CHE  & 2A      & 5 & FW & & General Chemistry \\
     & 2B      & 5 & WS & & General Chemistry \\
     & 2C      & 5 & FS & & General Chemistry \\
CHE  & 124A    & 3 & & & Inorganic Chemistry\\
\hline
\hline
PHY  & 115A   & 4 & & & \\
     & 115B   & 4 & & & \\
     & 112    & 4 & & & \\
     & 140A   & 4 & & & \\
     & 140B*   & 4 & & & \\
\hline
& Or: &&&& \\
\hline
CHE  & 110A   & 4 & & & Physical Chemistry: Intro to QM \\
     & 110B   & 4 & & & Physical Chemistry: Atoms and Molecules \\
     & 110C   & 4 & & & Physical Chemistry: Thermodynamics \\
     & 128A   & 4 & & & \\
     & 128B*   & 4 & & & \\
PHY  & 115A*  & 4 & & & \\
     & 112*   & 4 & & & \\
\hline
\hline
& 6 or more units: & & & & \\
PHY & 122A/B & 4 & & & \\
    & 117   & 4 & & & \\
    & 118   & 4 & & & \\
CHE & 105    & 4 & & & Analytical \& Physical Chemical Methods \\
    & 115    & 4 & & & Instrumental Analysis \\
    & 124L   & 2 & & & Laboratory Methods in Inorganic Chemistry \\
    & 129A   & 2 & & & Organic Chemistry Laboratory \\
    & 129B   & 2 & & & Organic Chemistry Laboratory \\
    & 129C   & 2 & & & Organic Chemistry Laboratory \\
\hline
\end{tabular}\\
\newpage

\letterpartb
        
\end{letter}

%----------------------------------------------------------------------------------
%  Computer Science, Chair Dipak Ghosal (ECS)
%----------------------------------------------------------------------------------
\begin{letter}{Department of Computer Science  \\ dghosal@ucdavis.edu}
\opening{Dear Chair Dipak Ghosal,}

\letterparta


{\bf No changes to ECS courses for Physics AB and BS majors:} There are no changes to Computer Science coursework to the Physics
majors (AB and BS).

{\bf Changes to ECS courses for Applied Physics majors:} The proposed physics curriculum includes changes to the Applied
Physics major with specialization in Computational Physics that impact
ECS coursework, as detailed in Table~1 (below).

In both the current and proposed major, students are required to take
ECS 36ABC and 122A.  In the proposed requirements, students choose at
least one elective from ECS 120,122B,132, and 171 (in the current
requirements, students choose from ECS 120, 122B, and 130).

It is worth noting that a major objective of the overall curriculum
update is to increase the role of computational physics in training
our students.  The specialization in Computational Physics benefits
from these improvements, including the new course PHY 40 (Introduction
to Computational Physics) and new one unit computational physics lab
courses (110L, 112L, and 115L).  Computational physics students are
expected to satisfy the prerequisites for 110L,112L, and 115L via ECS
36ABC, whereas other majors satisfy this requirement via the less
ambitious (computationally) PHY 45 (Computational Physics).

In summary, the changes to the Computational Physics specialization
are aimed at providing flexibility for students to study programming
and computational techniques outside of physics, and then bring those
skills back to the physics department to apply them specifically to
physics problems.

{\bf Table 1} Required ECS coursework in Computational Physics
\vskip 0.25cm
\noindent
*: recommended, $^\#$: not offered every year, $\parallel$: concurrently.\\
\begin{tabular}{|llllll|}
\hline
Course & & Units & Offered & Prereqs & Name \\
\hline
ECS & 36A  & 4 & FWS & & Programming and Problem Solving\\
    & 36B  & 4 & FWS & & Software Dev. and OOP in C++\\
    & 36C  & 4 & FWS & & Data Struct., Algorithms, and Prog.\\
    & 122A & 4 & FWS & & Algorithm Design \& Analysis\\
\hline
\hline
    & Choose one of & & & & \\
\hline
ECS & 120  & 4 & FWS & & Theory of Computation \\
    & 122B & 4 & WS  & & Algorithm Design \& Analysis \\
    & 132  & 4 & FWS & & Probability \& Stat Modeling \\
    & 171  & 4 & F   & & Machine Learning \\
\hline
\end{tabular}\\

{\bf Expected impact of these changes on the department:}  Based on past enrollment, we expect these changes to impact at most
five students per year per course:
\begin{quote}
{\bf ECS 36A/36B/36C/120/122A/122B/132/171}.
\end{quote}

Also, these majors include a Program Variance, which already allows
students to substitute related coursework with approval from the
department.  In fact, due to challenges with enrolling in popular ECS
offerings, most of our computational physics majors are already
substituting some of the proposed additional courses as electives on a
case-by-case basis.  {\bf We therefore expect the impact of these
  changes on CHE course enrollment to be quite minimal.}


\letterpartb
        
\end{letter}

%----------------------------------------------------------------------------------
%  Electrical \& Computer Engineering, Chair Andre Knoesen (EEC)
%----------------------------------------------------------------------------------
\begin{letter}{Department of Electrical and Computer Engineering \\ aknoesen@ucdavis.edu}
  
\opening{Dear Chair Andre Knoeson,}

\letterparta

{\bf No changes to EEC courses for Physics AB and BS majors:} There
are no changes to Electrical \& Computer Engineering coursework to the
Physics majors (AB and BS).

{\bf Changes to EEC courses and ENG 17 for Applied Physics majors:}  
The proposed physics curriculum includes changes to the Applied
Physics major with specialization in Physical Electronics that impact
ENG and EEC coursework, as detailed in Table~1 (below).

In both the current and proposed requirements, students are required
to take EEC 17 and 100.  In the proposed requirements, students choose
four elective courses from EEC 110A, 110B, 140A, 140B, 150, and 151.
(In the current requirements, 150 and 151 are not included as choices)

{\bf Expected impact of these changes on the Electrical \& Computer
  Engineering department:} Based on past enrollment, we expect these
changes to impact at most five students per year per course:
\begin{quote}
{\bf ENG 17, EEC 100/110A/110B/140A/140B/150/151}.
\end{quote}

These majors include a Program Variance, which already allows students
to substitute related coursework with approval from the department.
We have been accepting other EEC offerings (such as 150A, now 150) as
electives on a case-by-case basis when requested by students.  These
changes merely make this flexibility more explicit.  {\bf We therefore
  expect the impact of these changes on ENG and EEC course enrollment
  to be quite minimal.}

{\bf Table 1:} Required EEC coursework for Physical Electronics
\vskip 0.25cm
\noindent
Units:  8 units. *: recommended, $^\#$: not offered every year, $\parallel$: concurrently.\\
\begin{tabular}{|llllll|}
\hline
Course & & Units & Offered & Prereqs & Name \\
\hline
ENG & 17     & 4 & & & Circuits I\\
EEC & 100    & 4 & & ENG 17& Circuits II\\
\hline
\hline
    & Choose four of & & & & \\
\hline
EEC & 110A  & 4 & & EEC 100, 140A$\parallel$& Electronic Circuits I \\
    & 110B  & 4 & & EEC 110A& Electronic Circuits II\\
    & 140A  & 4 & & ENG 17 & Principles of Device Physics I\\
    & 140B  & 4 & & EEC 140B & Principles of Device Physics II\\
    & 150   & 4 & & EEC 100 & Intro. to Signals \& Systems \\
    & 151   & 4 & & EEC 100 & Digital Signals \& Systems \\
\hline
\end{tabular}\\

\letterpartb
        
\end{letter}

%----------------------------------------------------------------------------------
%  Materials Science \& Engineering, Chair Yayoi Takamura (EMS, ENG 45)
%----------------------------------------------------------------------------------
\begin{letter}{Department of Materials Science and Engineering \\ ytakamura@ucdavis.edu}

\opening{Dear Chair Yayoi Takamura,}

\letterparta

{\bf No changes to EMS and ENG courses for Physics AB and BS majors:}
There are no changes to Materials Science \& Engineering coursework
to the Physics majors (AB and BS).

{\bf Changes to EMS courses and ENG 45 for Applied Physics majors:}
The proposed physics curriculum includes changes to the Applied
Physics major with specializations in Materials Physics that impact
ENG and EMS coursework, as detailed in Table~1 (below).

In the current requirements, Materials Physics students are required
to take EMS 174 and 180, which implicitly requires them to take ENG 45
as well, pushing this degree past the unit cap on required coursework.
In the proposed requirements, students are required to take ENG 45,
and choose two electives from EMS 162, 160+164, 170, 172, 174, and
180.  Furthermore, they may satisfy their lab work requirements from
several choices within both PHY and EMS.

{\bf Table 1:} Required ENG and EMS coursework for Materials Physics
\vskip 0.25cm
\noindent
*: recommended, $^\#$: not offered every year, $\parallel$: concurrently.\\
\begin{tabular}{|llllll|}
\hline
Course & & Units & Offered & Prereqs & Name \\
\hline
ENG & 45     & 4 & FWS & & Properties of Materials \\
\hline
\hline
    & Choose two of & & & & \\
\hline
EMS & 162     & 4 & W & & Structure \& Characterization \\
    & 160+164 & 8 & F+W & & Thermo+Kinetics \\
    & 170     & 4 & F & & Sustainable Energy \\ 
    & 172     & 4 & S & & Smart Materials \\
    & 174     & 4 & S & & Mech. Behavior of Materials\\
    & 180     & 4 & F & & Materials in Eng. Design\\
\hline
\hline
    & Choose two of & & & & \\
\hline
PHY & 122A/B & 4 & & & \\
    & 117   & 4 & & & \\
    & 118   & 4 & & & \\
    & at most one of & & & & \\
EMS & 162L   & 3 & W & & Structure \& Characterization Lab\\
    & 170L   & 3 & S & & Sustainable Energy Lab \\
    & 172L   & 3 & F & & Smart Materials Lab \\
    & 174L   & 3 & S & & Mech. Behavior of Materials Lab \\
\hline
\end{tabular}\\

{\bf Expected impact of these changes on the Materials Science \& Engineering department:}
Based on past enrollment, we expect these changes to impact at most
five students per year per course:
\begin{quote}
{\bf ENG 45, EMS 160/162/162L/164/170/170L/172/172L/174/174L/180}.
\end{quote}
These majors include a Program Variance, which already allows students
to substitute related coursework with approval from the department.
For example, we would already under the current program accept EMS 170
as an elective if requested by a student.

These changes merely make this flexibility more explicit.  {\bf We
  therefore expect the impact of these changes on ENG and EMS course
  enrollment to be quite minimal.}


\letterpartb
        
\end{letter}

\end{document}


\section{Proposed BS with Specialization in Astrophysics}

The proposed required courses for the astrophysics specialization are presented in
Tables~\ref{tbl:prep-astro-applied} and \ref{tbl:depth-astro}.  To avoid
version conflicts, some information in Tables \ref{tbl:prep} and
\ref{tbl:depth} is not duplicated in these tables.

\captionof{table}{\label{tbl:prep-astro-applied}Preparatory Subject Matter:  Astrophysics, Applied Physics, and AB.}
\noindent
\vskip 0.25cm
\begin{center}
Units:  48-50. *: recommended.\\
\begin{tabular}{|lll|}
\hline
Course & & Units \\
\hline
MAT & 21A & 4 \\  
    & 21B & 4 \\ 
    & 21C & 4 \\ 
    & 21D & 4 \\
\hline
\hline
    & one of: & \\
\hline
MAT & 22A & 3 \\
    & 27A & 3 \\
    & 67  & 4 \\
\hline
\hline
MAT & 22B & 3 \\ 
    & or: & \\
MAT & 27B & 3 \\ 
\hline
\hline
PHY & 9A & 5 \\  
    & 9B & 5 \\  
    & 9C & 5 \\  
    & 9D & 4 \\  
\hline
&or&\\
\hline
PHY & 9HA & 5 \\  
    & 9HB & 5 \\  
    & 9HC & 5 \\  
    & 9HD & 5 \\  
\hline
\hline
PHY & 40   & 3 \\
    & 80   & 4 \\ 
    & 185* & 1 \\ 
    & 190* & 1 \\ 
\hline
\end{tabular}
\end{center}

\newpage
\captionof{table}{\label{tbl:depth-astro}Depth Subject Matter: Astrophysics}
\noindent
\vskip 0.25cm
Units:  60-64. *: recommended, $^\#$: not offered every year, $\parallel$: concurrently.\\
\begin{tabular}{|llllll|}
\hline
Course & & Units & Offered & Prereqs & Name \\
\hline
AST & 25* & 4 & & & \\ 
PHY & 104A & 4 & & & \\ 
    & 105A & 4 & & & \\
    & 108  & 3 & S & 9D/9HD & Optics \\  
    & 108L & 1 & S & $\parallel$108 & Optics Laboratory \\  
    & 110A & 4 & & & \\
    & 110B & 4 & & & \\
    & 112  & 4 & & & \\    
    & 115A & 4 & & & \\
    & 115B & 4 & & & \\
\hline
\hline
PHY & 157 & 4 & S$^\#$ & (Same as 122A/B) & Astronomy Instrumentation and \\  
    &     &   &     & & Data Analysis Lab\\  
\hline
    & or & & & & \\
\hline
PHY & 122A/B & 4 & & & Advanced Physics Laboratory \\  
\hline
\hline
 & Choose four of & & & \\
\hline
PHY & 151 & 4 & F$^\#$ & $\parallel$40,$\parallel$104A & Stellar Structure and Evolution \\ 
    & 152 & 4 & F$^\#$ & $\parallel$40,$\parallel$104A & Galactic Structure and \\
    &     &   &     &           & the Interstellar Medium\\
    & 153 & 4 & W$^\#$ & 40,104A,$\parallel$105A & Extragalactic Astrophysics\\  
    & 156 & 4 & W$^\#$ & 40,104A,$\parallel$105A & Introduction to Cosmology\\ 
    & 158  & 4 & S$^\#$ & 40,104A,105A & Galaxy Formation \\ 
\hline
 & Any two of & & & & : \\
\hline 
PHY & 105B & 4 &  &  & \\ 
    & 117 & 4 &  &  & \\  
    & 118 & 4 &  &  & \\  
    & 129A & 4 &  &  & \\  
    & 130A & 4 &  &  & \\  
    & 130B & 4 &  &  & \\  
    & 150  & 4 &  &  & Special Topics\\  
    & 154  & 4 & S$^\#$ & 40,105B,110B,115A & Astro. Appl. of Phys. \\
    & 155  & 4 & W   & 104A,105B,110A & General Relativity \\ 
GEL & 163  & 4 &  &  & Planetary Geology and Geophysics\\ 
 & At most one of &  &  &  & \\ 
PHY & 194HAB & 8 &  &  & Special Study for Honors Students\\ 
    & 195  & 5 &  &  & Senior Thesis\\
    & 198  & 3+ &  &  & Directed Group Study\\ 
    & 199  & 3+ &  &  & Special Study for Adv. Undergrads.\\ 
\hline
\end{tabular}\\
\vskip 0.25cm
\noindent
PHY 108 has alternative prerequisites (PHY 7C and 21D)
intended for non-majors.  PHY 150 must be an astro topic and requires
prior department approval.\\
\noindent
{\bf Total Units:} 108-114 \\


\newpage
\section{Proposed Applied Physics Majors}

The applied physics majors (excluding chemical physics) require the
preparatory courses in Table~\ref{tbl:prep-astro-applied} and the
depth courses listed in Table~\ref{tbl:depth-applied}.  Together,
these required courses account for 72-74 units, leaving 38 units for
concentration courses as detailed in the following tables.  This
proposal includes modifications to all of the applied physics majors:
\begin{itemize}

  \item The Computational Physics major does not require PHY 45 and ECS
  36B will satisfy the prerequisite for the computing labs.  The
  elective choices from the Computer Science and Engineering
  department have been extended.  See Table~\ref{tbl:computational}.

\item The Physical Electronics major does not include 117 and 118 as
  part of the upper division lab requirement, only 122A/B is required.
  The major instead includes a large amount of electronics from the
  Electrical and Computer Engineering department.  Due to unit limits,
  only one computational lab is required.  See
  Table~\ref{tbl:electronics}.

\item The Materials Science major is renamed Materials Physics, and
  the elective choices have been extended.  Students may now use a
  Material Science and Engineering lab course as one of their required
  upper division lab courses.  See Table~\ref{tbl:materials}.

\item In the past, the Atmospheric Physics major required the Physical
  and Chemical Oceanography course (GEL 150A), which now has extensive
  prerequisites that would exceed the unit limit.  The Oceanography
  thrust of this course has been relocated to the electives, including
  the more easily obtained (GEL 116N).  The list of Atmospheric
  Physics electives has been extended.  See Table~\ref{tbl:atmospheric}.

\item The Physical and Chemical Oceanography course (GEL 150A) is not
  offered every year and now has extensive prerequisites.  Therefore,
  the Physical Oceanography major might not be achievable in two years
  for students that have not already satisfied these prerequisites.
  The prerequisites for reaching GEL 150A are now included in the
  concentration courses.  See Table~\ref{tbl:oceanography}.
  
\item The Chemical Physics major requires the courses listed in
  Tables~\ref{tbl:prep-astro-applied} and~\ref{tbl:chemical}.  
  Both chemistry and physics have significant lower division
  coursework as prerequisite for upper division coursework.  The
  additional 15 units of lower division general chemistry course
  work required by this major limits the number of upper division
  courses that can be required.  This major features two paths, one
  featuring PHY 112, 115AB, 140AB and the other featuring CHE 110ABC
  128AB.  The latter choice relies on the treatment of quantum
  mechanics and thermodynamics provided by our colleagues in the
  chemistry department.  This major requires 6 units of upper division
  lab work in chemistry or physics.
\end{itemize}
Some of the majors might benefit from further revisions, which we plan
to pursue once we have experience with the significant changes already
proposed here.  For example, it would be helpful if PHY 115AB were
listed as an alternative prerequisite for CHE 110B.

\captionof{table}{\label{tbl:depth-applied} Depth subject matter for
  all applied physics majors excluding chemical physics.}
\noindent
\vskip 0.25cm
\begin{center}
\begin{tabular}{|lll|}
\hline
Course & & Units \\
\hline
PHY & 104A   & 4 \\  
    & 105A   & 4 \\ 
    & 110A   & 4 \\ 
    & 110B   & 4 \\
    & 112    & 4 \\
    & 115A   & 4 \\ 
\hline
\end{tabular}
\end{center}

\newpage
\captionof{table}{\label{tbl:computational}Computational Physics}
\vskip 0.25cm
\noindent
{\bf Additional Preparatory Courses:  }\\
Units:  12 units. *: recommended, $^\#$: not offered every year, $\parallel$: concurrently.\\
\begin{tabular}{|llllll|}
\hline
Course & & Units & Offered & Prereqs & Name \\
\hline
ECS & 36A  & 4 & FWS & & Programming and Problem Solving\\
    & 36B  & 4 & FWS & & Software Dev. and OOP in C++\\
    & 36C  & 4 & FWS & & Data Structures, Algorithms, and Programming\\
\hline
\end{tabular}\\
\vskip 0.25cm
\noindent
{\bf Concentration Courses:  }\\
Units:  18 units. *: recommended, $^\#$: not offered every year, $\parallel$: concurrently.\\
\begin{tabular}{|llllll|}
\hline
Course & & Units & Offered & Prereqs & Name \\
\hline
ECS & 122A & 4 & FWS & & Algorithm Design \& Analysis\\
\hline
\hline
    & Choose two of & & & & \\
\hline
PHY    & 110L & 1 & & & \\
    & 115L & 1 & & & \\
    & 112L & 1 & & & \\
\hline
\hline
    & Choose one of & & & & \\
\hline
ECS & 120  & 4 & FWS & & Theory of Computation \\
    & 122B & 4 & WS  & & Algorithm Design \& Analysis \\
    & 132  & 4 & FWS & & Probability \& Stat Modeling \\
    & 171  & 4 & F   & & Machine Learning \\
\hline
\hline
    & Choose two of & & & & \\
\hline
PHY & 122A/B & 4 & & & \\
    & 117   & 4 & & & \\
    & 118   & 4 & & & \\
\hline
\end{tabular}\\
\vskip 0.25cm
\noindent
{\bf Additional Electives: (8 units)} Two courses for a total of at least 8 units of additional upper division coursework from ECS, MAT, or PHY.\\
{\bf Total Units:} 110-112\\



\newpage
\captionof{table}{\label{tbl:electronics}Physical Electronics }
\vskip 0.25cm
\noindent
{\bf Additional Preparatory Courses:  }\\
Units:  8 units. *: recommended, $^\#$: not offered every year, $\parallel$: concurrently.\\
\begin{tabular}{|llllll|}
\hline
Course & & Units & Offered & Prereqs & Name \\
\hline
ENG & 17     & 4 & FWS & & Circuits I\\
PHY & 45     & 4 & & & \\
\hline
\end{tabular}\\
\vskip 0.25cm
\noindent
{\bf Concentration Courses:  }\\
Units:  30 units. *: recommended, $^\#$: not offered every year, $\parallel$: concurrently.\\
\begin{tabular}{|llllll|}
\hline
Course & & Units & Offered & Prereqs & Name \\
\hline
EEC & 100    & 4 & FW & & Circuits II\\
PHY & 115B*  & 4 & & & \\
    & 140A   & 4 & & & \\
    & 140B*  & 4 & & & \\
    & 122A/B & 4 & & & \\
\hline
\hline
    & Choose two of & & & & \\
\hline
PHY & 110L & 1 & & & \\
    & 115L & 1 & & & \\
    & 112L & 1 & & & \\
\hline
\hline
    & Choose four of & & & & \\
\hline
EEC & 110A  & 4 & WS & & Electronic Circuits I \\
    & 110B  & 4 & S  & & Electronic Circuits II\\
    & 140A  & 4 & FW & & Principles of Device Physics I\\
    & 140B  & 4 & S  & & Principles of Device Physics II\\
    & 150A  & 4 & WS & & Intro. to Signals \& Systems I\\
    & 150B  & 4 & F  & & Intro. to Signals \& Systems II\\
\hline
\end{tabular}\\
\vskip 0.25cm
\noindent
{\bf Total Units:} 110-112\\

\newpage
\captionof{table}{\label{tbl:materials}Materials Physics}
\vskip 0.25cm
\noindent
{\bf Additional Preparatory Courses:  }\\
Units:  8 units. *: recommended, $^\#$: not offered every year, $\parallel$: concurrently.\\
\begin{tabular}{|llllll|}
\hline
Course & & Units & Offered & Prereqs & Name \\
\hline
PHY & 45     & 4 & & & \\
ENG & 45     & 4 & FWS & & Properties of Materials \\
\hline
\end{tabular}\\
\vskip 0.25cm
\noindent
{\bf Concentration Courses:  }\\
Units:  30-34 units. *: recommended, $^\#$: not offered every year, $\parallel$: concurrently.\\
\begin{tabular}{|llllll|}
\hline
Course & & Units & Offered & Prereqs & Name \\
\hline
PHY & 115B   & 4 & & & \\
    & 140A   & 4 & & & \\
    & 140B   & 4 & & & \\
    & 110L & 1 & & & \\
    & 115L & 1 & & & \\
    & 112L & 1 & & & \\
\hline
\hline
    & Choose two of & & & & \\
\hline
EMS & 162     & 4 & W & & Structure \& Characterization \\
    & 160+164 & 7 & F+W & & Thermo+Kinetics \\
    & 170     & 4 & S & & Sustainable Energy \\ 
    & 172     & 4 & F & & Smart Materials \\
    & 174     & 4 & S & & Mech. Behavior of Materials\\
    & 180     & 4 & F & & Materials in Eng. Design\\
\hline
\hline
    & Choose two of & & & & \\
\hline
PHY & 122A/B & 4 & & & \\
    & 117   & 4 & & & \\
    & 118   & 4 & & & \\
    & at most one of & & & & \\
EMS & 162L   & 3 & W & & Structure \& Characterization Lab\\
    & 170L   & 3 & S & & Sustainable Energy Lab \\
    & 172L   & 3 & F & & Smart Materials Lab \\
    & 174L   & 3 & S & & Mech. Behavior of Materials Lab \\
\hline
\end{tabular}\\
\noindent
{\bf Total Units:} 110-116 \\

\newpage
\captionof{table}{\label{tbl:atmospheric}Atmospheric Physics}
\vskip 0.25cm
\noindent
{\bf Additional Preparatory Courses:  }\\
Units: 11 units. *: recommended, $^\#$: not offered every year, $\parallel$: concurrently.\\
\begin{tabular}{|llllll|}
\hline
Course & & Units & Offered & Prereqs & Name \\
\hline
PHY & 45     & 4 &     & & \\
GEL & 50     & 3 & FWS & & Physical Geology \\
ATM & 60     & 4 & F   & & Intro. to Atmospheric Sci. \\
\hline
\end{tabular}\\
\vskip 0.25cm
\noindent
{\bf Concentration Courses:  }\\
Units:  27-29 units. *: recommended, $^\#$: not offered every year, $\parallel$: concurrently.\\
\begin{tabular}{|llllll|}
\hline
Course & & Units & Offered & Prereqs & Name \\
\hline
ATM & 120    & 4 & F   & & Atmos. Thermo. and Cloud Physics \\
    & 121A   & 4 & W   & & Atmospheric Dynamics \\
    & 121B   & 4 & S   & & Atmospheric Dynamics \\
\hline
\hline
    & Choose one of & & & & \\
\hline
PHY & 110L   & 1 & & & \\
    & 115L   & 1 & & & \\
    & 112L   & 1 & & & \\
\hline
\hline
    & Choose two of & & & & \\
\hline
PHY & 122A/B & 4 & & & \\
    & 117   & 4 & & & \\
    & 118   & 4 & & & \\
\hline
\hline
    & Choose two of & & & & \\
\hline
PHY  & 105B   & 4 & & & \\
     & 105C   & 4 & $^\#$  &             & Continuum Mechanics\\
GEL  & 116N   & 3 & S  & GEL 50      & Oceanography\\
     & 150A   & 4 & S$^\#$ & GEL 116N,55 & Physical and Chemical Oceanography\\
     & 150B   & 3 & W  & GEL 50      & Geological Oceanography\\
ATM  & 124    & 3 & F  & ATM 60      & Meteorological Instruments \& Observations \\
     & 128    & 4 & W  & ATM 60      & Radiation and Satellite Meteorology \\
     & 158    & 4 & S  & ATM 121A    & Boundary-Layer Meteorology \\
\hline
\end{tabular}\\
\vskip 0.25cm
\noindent
{\bf Total Units:} 110-114 \\


\newpage
\captionof{table}{\label{tbl:oceanography}Physical Oceanography}
\vskip 0.25cm
\noindent
{\bf Additional Preparatory Courses:  }\\
Units:  15 units. *: recommended, $^\#$: not offered every year, $\parallel$: concurrently.\\
\begin{tabular}{|llllll|}
\hline
Course & & Units & Offered & Prereqs & Name \\
\hline
CHE  & 2A     & 5 & FW  & & \\
PHY  & 45     & 4 &     & & \\
GEL  & 50     & 3 & FWS & & Physical Geology \\
     & 55     & 3 & F      & CHE 2A & Intro. to Geochemistry\\
ATM  & 60*    & 4 & F   & & Intro. to Atmospheric Sci. \\
\hline
\end{tabular}\\
\vskip 0.25cm
\noindent
{\bf Concentration Courses:  }\\
Units:  23-25 units. *: recommended, $^\#$: not offered every year, $\parallel$: concurrently.\\
\begin{tabular}{|llllll|}
\hline
Course & & Units & Offered & Prereqs & Name \\
\hline
GEL  & 116N   & 3 & S      & GEL 50          & Oceanography\\
     & 150A   & 4 & S$^\#$ & GEL 116N,55 & Physical and Chemical Oceanography\\
\hline
    & Choose two of & & & & \\
\hline
PHY  & 110L & 1 & & & \\
     & 112L & 1 & & & \\
     & 115L & 1 & & & \\
\hline
\hline
    & Choose two of & & & & \\
\hline
PHY & 122A/B & 4 & & & \\
    & 117   & 4 & & & \\
    & 118   & 4 & & & \\
\hline
\hline
    & Choose two of & & & & \\
\hline
PHY  & 105B  & 4 & & & \\
     & 105C  & 4 & S$^\#$ & & Continuum Mechanics \\
ATM  & 115   & 3 & S   & ATM 60 & Hydroclimatology \\
     & 116   & 3 & F   & & Modern Climate Change \\
     & 120   & 4 & F   & ATM 60$\parallel$ & Atmos. Thermo. and Cloud Physics \\
GEL  & 150B  & 3 & W & GEL 50 & Geological Oceanography\\
\hline
\end{tabular}\\
\vskip 0.25cm
\noindent
{\bf Total Units:} 110-114\\
{\bf Note: } For transfer students arriving in their junior year, this
degree may not be possible to complete in two years.  Students should
seek instructor permission to take GEL 150A alongside 116N prerequisite if
it is offered in their junior year.\\

\newpage
\captionof{table}{\label{tbl:chemical}Chemical Physics}
\vskip 0.25cm
\noindent
{\bf Additional Preparatory Courses:  }\\
Units:  19 units. *: recommended, $^\#$: not offered every year, $\parallel$: concurrently.\\
\begin{tabular}{|llllll|}
\hline
Course & & Units & Offered & Prereqs & Name \\
\hline
\hline
CHE  & 2A      & 5 & FW & & General Chemistry \\
     & 2B      & 5 & WS & & General Chemistry \\
     & 2C      & 5 & FS & & General Chemistry \\
PHY  & 45      & 4 & & & \\
\hline
\end{tabular}\\
\noindent
{\bf Concentration Courses:  }\\
Units:  43-45 units. *: recommended, $^\#$: not offered every year, $\parallel$: concurrently.\\
\begin{tabular}{|llllll|}
\hline
Course & & Units & Offered & Prereqs & Name \\
\hline
PHY & 104A     & 4 & & & \\  
    & 105A     & 4 & & & \\ 
    & 110A     & 4 & & & \\ 
    & 110B     & 4 & & &\\    
CHE & 124A     & 3 & & & Inorganic Chemistry\\
\hline
\hline
PHY  & 115A   & 4 & & & \\
     & 115B   & 4 & & & \\
     & 112    & 4 & & & \\
     & 140A   & 4 & & & \\
     & 140B*   & 4 & & & \\
\hline
& Or: &&&& \\
\hline
CHE  & 110A   & 4 & & & Physical Chemistry: Intro to QM \\
     & 110B   & 4 & & & Physical Chemistry: Atoms and Molecules \\
     & 110C   & 4 & & & Physical Chemistry: Thermodynamics \\
     & 128A   & 4 & & & \\
     & 128B*   & 4 & & & \\
PHY  & 115A*  & 4 & & & \\
     & 112*   & 4 & & & \\
\hline
\hline
& Choose two of: &&&& \\
\hline
PHY  & 110L  & 1 & & & \\
     & 112L  & 1 & & & \\
     & 115L  & 1 & & & \\
\hline
& 6 or more units: & & & & \\
PHY & 122A/B & 4 & & & \\
    & 117   & 4 & & & \\
    & 118   & 4 & & & \\
CHE & 105    & 4 & & & Analytical \& Physical Chemical Methods \\
    & 115    & 4 & & & Instrumental Analysis \\
    & 124L   & 2 & & & Laboratory Methods in Inorganic Chemistry \\
    & 129A   & 2 & & & Organic Chemistry Laboratory \\
    & 129B   & 2 & & & Organic Chemistry Laboratory \\
    & 129C   & 2 & & & Organic Chemistry Laboratory \\
\hline
\end{tabular}\\
\noindent
{\bf Total Units:} 110-114 \\

\newpage
\captionof{table}{\label{tbl:geophysics} Geological Physics}
\vskip 0.25cm
\noindent
{\bf Additional Preparatory Courses:  }\\
Units: 13 units. *: recommended, $^\#$: not offered every year, $\parallel$: concurrently.\\
\begin{tabular}{|llllll|}
\hline
Course & & Units & Offered & Prereqs & Name \\
\hline
PHY & 45     & 4 &     & & \\
GEL & 50     & 3 & FWS & & Physical Geology \\
    & 50L    & 2 &     & & Physical Geology Laboratory\\
ATM & 60     & 4 & F   & & Intro. to Atmospheric Sci. \\
\hline
\end{tabular}\\
\vskip 0.25cm
\noindent
{\bf Concentration Courses:  }\\
Units:  25-30 units. *: recommended, $^\#$: not offered every year, $\parallel$: concurrently.\\
\begin{tabular}{|llllll|}
\hline
Course & & Units & Offered & Prereqs & Name \\
\hline
    & Choose two of & & & & \\
\hline
PHY & 110L   & 1 & & & \\
    & 115L   & 1 & & & \\
    & 112L   & 1 & & & \\
\hline
\hline
    & Choose two of & & & & \\
\hline
PHY & 122A/B & 4 & & & \\
    & 117   & 4 & & & \\
    & 118   & 4 & & & \\
\hline
\hline
    & Choose five of & & & & \\
\hline
PHY  & 105B   & 4 & &  & \\
     & 105C   & 4 & &  & Continuum Mechanics\\
ATM  & 120    & 4 & F  & ATM 60    & Atmospheric Thermodynamics \& Cloud Physics \\
     & 121A   & 4 & F  & ATM 120   & Atmospheric Dynamics \\
     & 121B   & 4 & W  & ATM 121A  & Atmospheric Dynamics \\
GEL  & 116N   & 3 &    & & Oceanography\\
     & 146    & 3 &    & & Radiogenic Isotope Geochemistry \& Cosmochemistry \\
     & 150A   & 4 & S$^\#$ & GEL 116N,55 & Physical and Chemical Oceanography\\
     & 150B   & 3 & W  & GEL 50      & Geological Oceanography\\
\hline
     & incl at least two:  & & & & \\
\hline
GEL  & 161    & 3 & S$^\#$ & GEL 50 & Geophysical Field Methods \\
     & 162    & 3 & W$^\#$ & GEL 50 & Geophysics of the Solid Earth \\
     & 163    & 3 & F$^\#$ & GEL 50 & Planetary Geology \\
\hline
\end{tabular}\\
\vskip 0.25cm
\noindent
{\bf Total Units:} 110-117 \\



\newpage
\section{Proposed AB Physics Major}
\label{sec:ab}

The proposed requirements for the AB Physics major are listed in
Tables~\ref{tbl:prep-astro-applied} and \ref{tbl:ab-depth}.  To avoid
version conflicts, some information in Tables \ref{tbl:prep} and
\ref{tbl:depth} is not duplicated in these tables.

\captionof{table}{\label{tbl:ab-depth}Depth Subject Matter}
\noindent
\vskip 0.25cm
\begin{center}
Units:  32. *: recommended, $\parallel$: concurrently.\\
\begin{tabular}{|lll|}
\hline
Course & & Units \\
\hline
PHY & 104A & 4 \\
    & 105A & 4 \\
    & 110A & 4 \\
    & 110B & 4 \\
    & 112  & 4 \\
    & 115A & 4 \\
\hline
\hline
PHY & 122A/B & 4 \\
\hline
    & At least one of & \\ 
\hline
PHY & 129A & 4 \\
    & 130A & 4 \\
    & 140A & 4 \\
    & 151 & 4 \\
    & 152 & 4 \\
    & 153 & 4 \\
    & 156 & 4 \\
    & 158 & 4 \\
\hline
\end{tabular}\\ \vskip 0.25cm
\end{center}
\noindent
{\bf Foreign Language:} The AB degree requires proficiency in a language other than English (15 units)\\
\noindent
{\bf Total Units:} 80-82 units, excluding language requirement.




\end{document}

